% GENERATED FILE. Edit canonical lessons, then run npm run generate:books.
% canonical-source-hash: fnv1a64:3c3bf87cc968fdd2
% canonical-lessons: RU-C150-redko, RU-C150-kazhdyy, RU-C150-snachala, RU-C150-nakonets, RU-C150-yeshchyo

\chapter{Rarely, Every, At first, At last, Still}
\label{ch:ru-rarely-every-at-first-at-last-still}
\hlchaptermodality{150}

\begin{chapteropening}
\textbf{By the end of this chapter:} \emph{I can say rarely, every, at first, at last and still.}

Complete the last lesson of chapter 150: I can say rarely, every, at first, at last and still.
\end{chapteropening}

\section[\texorpdfstring{\ru{редко} (rédko)}{редко (rédko)}]{\ru{редко} (rédko) — rarely}

\label{lesson:RU-C150-redko}



\begin{quote}
Before the new one: say the Russian for dawn, then the Russian for sunset.
\end{quote}

\subsection*{You'll want to know: \ru{редко}}
\textbf{\ru{редко}} (\emph{rédko}) — \textquotedblleft{}rarely\textquotedblright{}.

\begin{grammarlens}[title={What you've built}]
One more word for this chapter.
\end{grammarlens}

\subsection*{Your turn}
\begin{itemize}
\raggedright
  \item \emph{Say it:} \emph{rédko}
  \item \emph{Say it:} \emph{rédko}, once more
  \item \emph{Say it:} say \emph{rédko}
  \item \emph{Recall:} say the Russian for dawn, then the Russian for sunset, then say \emph{rédko} again
\end{itemize}

\begin{tcolorbox}[breakable,colback=teal!4,colframe=teal!35!black,title={Before you move on}]
What does \ru{редко} mean? (\textquotedblleft{}Rarely\textquotedblright{}.) Say it once more.
\end{tcolorbox}
\section[\texorpdfstring{\ru{каждый} (kázhdyy)}{каждый (kázhdyy)}]{\ru{каждый} (kázhdyy) — every, each}

\label{lesson:RU-C150-kazhdyy}



\begin{quote}
Before the new one: say the Russian for sunset, then the Russian for rarely.
\end{quote}

\subsection*{You'll want to know: \ru{каждый}}
\textbf{\ru{каждый}} (\emph{kázhdyy}) — \textquotedblleft{}every, each\textquotedblright{}.

\begin{grammarlens}[title={What you've built}]
One more word for this chapter.
\end{grammarlens}

\subsection*{Your turn}
\begin{itemize}
\raggedright
  \item \emph{Say it:} \emph{kázhdyy}
  \item \emph{Say it:} \emph{kázhdyy}, once more
  \item \emph{Say it:} say \emph{kázhdyy}
  \item \emph{Recall:} say the Russian for sunset, then the Russian for rarely, then say \emph{kázhdyy} again
\end{itemize}

\begin{tcolorbox}[breakable,colback=teal!4,colframe=teal!35!black,title={Before you move on}]
What does \ru{каждый} mean? (\textquotedblleft{}Every, each\textquotedblright{}.) Say it once more.
\end{tcolorbox}
\section[\texorpdfstring{\ru{сначала} (snachála)}{сначала (snachála)}]{\ru{сначала} (snachála) — at first}

\label{lesson:RU-C150-snachala}



\begin{quote}
Before the new one: say the Russian for rarely, then the Russian for every.
\end{quote}

\subsection*{You'll want to know: \ru{сначала}}
\textbf{\ru{сначала}} (\emph{snachála}) — \textquotedblleft{}at first\textquotedblright{}.

\begin{grammarlens}[title={What you've built}]
One more word for this chapter.
\end{grammarlens}

\subsection*{Your turn}
\begin{itemize}
\raggedright
  \item \emph{Say it:} \emph{snachála}
  \item \emph{Say it:} \emph{snachála}, once more
  \item \emph{Say it:} say \emph{snachála}
  \item \emph{Recall:} say the Russian for rarely, then the Russian for every, then say \emph{snachála} again
\end{itemize}

\begin{tcolorbox}[breakable,colback=teal!4,colframe=teal!35!black,title={Before you move on}]
What does \ru{сначала} mean? (\textquotedblleft{}At first\textquotedblright{}.) Say it once more.
\end{tcolorbox}
\section[\texorpdfstring{\ru{наконец} (nakonéts)}{наконец (nakonéts)}]{\ru{наконец} (nakonéts) — at last, finally}

\label{lesson:RU-C150-nakonets}



\begin{quote}
Before the new one: say the Russian for every, then the Russian for at first.
\end{quote}

\subsection*{You'll want to know: \ru{наконец}}
\textbf{\ru{наконец}} (\emph{nakonéts}) — \textquotedblleft{}at last, finally\textquotedblright{}.

\begin{grammarlens}[title={What you've built}]
One more word for this chapter.
\end{grammarlens}

\subsection*{Your turn}
\begin{itemize}
\raggedright
  \item \emph{Say it:} \emph{nakonéts}
  \item \emph{Say it:} \emph{nakonéts}, once more
  \item \emph{Say it:} say \emph{nakonéts}
  \item \emph{Recall:} say the Russian for every, then the Russian for at first, then say \emph{nakonéts} again
\end{itemize}

\begin{tcolorbox}[breakable,colback=teal!4,colframe=teal!35!black,title={Before you move on}]
What does \ru{наконец} mean? (\textquotedblleft{}At last, finally\textquotedblright{}.) Say it once more.
\end{tcolorbox}
\section[\texorpdfstring{\ru{ещё} (yeshchyó)}{ещё (yeshchyó)}]{\ru{ещё} (yeshchyó) — still; more}

\label{lesson:RU-C150-yeshchyo}



\begin{quote}
Before the new one: say the Russian for at first, then the Russian for at last.
\end{quote}

\subsection*{You'll want to know: \ru{ещё}}
\textbf{\ru{ещё}} (\emph{yeshchyó}) — \textquotedblleft{}still; more\textquotedblright{}.

\begin{grammarlens}[title={What you've built}]
One more word for this chapter.
\end{grammarlens}

\subsection*{Your turn}
\begin{itemize}
\raggedright
  \item \emph{Say it:} \emph{yeshchyó}
  \item \emph{Say it:} \emph{yeshchyó}, once more
  \item \emph{Say it:} say \emph{yeshchyó}
  \item \emph{Recall:} say the Russian for at first, then the Russian for at last, then say \emph{yeshchyó} again
\end{itemize}

\begin{tcolorbox}[breakable,colback=teal!4,colframe=teal!35!black,title={Before you move on}]
What does \ru{ещё} mean? (\textquotedblleft{}Still; more\textquotedblright{}.) Say it once more.
\end{tcolorbox}
